% !TEX root = ../main.tex
\lecture{4}{2026-10-08}{Nutzenmaximierung}

\subsection{Indifferenz als Funktion}
Die einzelnen Höhenlinien werden als Indifferenzkurven behandelt. Dabei sind Waren in einer bestimmten Menge substituierbar, ohne dass Nutzen verloren wird.

\subsubsection{Eigenschaften}
\begin{itemize}
    \item Konvexität
    \item Absolutewerte der Nutzen haben keine Bedeutung, nur relativ zu einander
    \item 2 Indifferenzkurven/Höhenlinien schneiden sich nicht, \textbf{keine gemeinsamen Punkte}
          \begin{itemize}
              \item Weil sonst keine Rangordnung möglich
          \end{itemize}
\end{itemize}

\subsubsection{Funktion}
Es wird ein Nutzenwert ausgewählt auf dessen Höhenlinie alle Nutzensgleiche Güterkombinationen befinden.
\begin{equation}
    U(W,F) = U_0 \rightarrow F(W) = \frac{U_0}{W}
\end{equation}

\pagebreak
\subsubsection{Cobb-Douglas Funktion}

\begin{definition}
    [Cobb-Douglas]
    \begin{equation}
        U(X,Y) = cX^\alpha Y^\beta
    \end{equation}

    \begin{equation}
        Y= \left[\frac{U_0}{cX}\right]^\frac{1}{\beta}
    \end{equation}
\end{definition}

\subsection{Grenzrate der Substitution (Marginal Rate of Substitution)}

\begin{definition}
    [GRS] Wie Verändert sich das Mengenverhältnis zweier Waren auf einer gleichen Indifferenzkurve zu jedem Punkt. Ableitung der Indifferenzfunktion.
    \begin{equation}
        F'(W) = \frac{U_0}{W}
    \end{equation}
\end{definition}

\subsubsection{Optimum}
Das Optimum befindet sich beim Totalen Differenzial der Nutzenfunktion beim Wert 0, sodass:\footnote{aus Multivariable Calculus}
\begin{equation}
    GRS(X^*,Y^*) = \frac{dU_X}{dU_Y}
\end{equation}

\subsection{Angewandte Mikroökonomie II}

\subsubsection{Fremdbezogene Präferenzen}

Gegeben sei eine Modellökonomie worin Akteure: A (egoistisch) und B (altruistisch), folgende Nutzenfunktionen haben. Zur Vereinfachung besteht jeder Güterbündel aus einer einzigen Ware.

\begin{equation}
    U_A(X_A) = X_A 
\end{equation}
\footnote{diskrete senkrechte Geraden}
\begin{equation}
    U_B(X_B, X_A)= {X_A}^{0.8}X_B^{0.2} 
\end{equation}
\footnote{diskrete Indifferenzkurven}
\subsubsection{Ungleichheitsaversion}


\begin{equation}
    U_B(X_A,X_B) = X_B - \alpha max\left\{X_A-X_B,0\right\} - \beta max \left\{ X_B - X_A, 0\right\}
\end{equation}

wobei angenommen wird, dass $0 < \beta < 1$ und $\beta \le \alpha$